% Q15_problems.tex — Practice II-C + Lab 2, written problems (QPLAN §25.6, §25.7).
% B1 (the pre-registration) and B2 (the lab report) are Lab 2's; the three paper cards
% (04.6, 04.7, 03.8) follow as B3-B5, 15 minutes each, before B1.
\documentclass[11pt]{article}
% preamble.tex — shared preamble for the dafny-verus-learn-0 written track.
%
% Usage (every <ID>_problems.tex, <ID>_answers.tex and solutions/<ID>_solutions.tex):
%
%   \documentclass[11pt]{article}
%   % preamble.tex — shared preamble for the dafny-verus-learn-0 written track.
%
% Usage (every <ID>_problems.tex, <ID>_answers.tex and solutions/<ID>_solutions.tex):
%
%   \documentclass[11pt]{article}
%   % preamble.tex — shared preamble for the dafny-verus-learn-0 written track.
%
% Usage (every <ID>_problems.tex, <ID>_answers.tex and solutions/<ID>_solutions.tex):
%
%   \documentclass[11pt]{article}
%   \input{preamble}          % the Makefile always runs latexmk from paper/
%   \begin{document}
%   ...
%   \end{document}
%
% Engine: pdflatex. Dafny and Verus sources are mostly ASCII; the `literate`
% tables below map the common Unicode symbols to math glyphs inside listings.
% For other Unicode, extend a table, or build that file with lualatex
% (`latexmk -lualatex <file>`) after replacing the inputenc/fontenc lines with
% \usepackage{fontspec} and a monospace font with good coverage.
%
% Contents: program-logic macros (Hoare triples, wp, VCs, substitution,
% semantics brackets, typing judgments), inference rules (mathpartir),
% theorem and problem environments, the answer environment, the grader
% margin-note macro, and listings languages `dafny` and `verus` (Python uses
% listings' built-in `Python`), with inline forms \dfy{...}, \vrs{...}, \py{...}.
\usepackage[utf8]{inputenc}
\usepackage[T1]{fontenc}

\usepackage{amsmath}
\usepackage{amssymb}
\usepackage{amsthm}
\usepackage{stmaryrd}         % \llbracket \rrbracket for semantics brackets
\usepackage{mathpartir}       % inference rules: \inferrule, the mathpar environment
\usepackage{listings}
\usepackage{xcolor}
\usepackage[margin=1.1in]{geometry}

% ---------------------------------------------------------------------------
% Program logic.
%   \hoare{P}{c}{Q}       {P} c {Q}           (partial correctness)
%   \thoare{P}{c}{Q}      [P] c [Q]           (total correctness)
%   \provable{P}{c}{Q}    |- {P} c {Q}        (derivable in the proof system)
%   \WP{c}{Q}             wp(c, Q)            (\wp is LaTeX's Weierstrass p, so
%   \WLP{c}{Q}            wlp(c, Q)            the macro names are upper case)
%   \VC{c}{Q}             VC(c, Q)
%   \subst{Q}{x}{e}       Q[e/x]              (e replaces x in Q)
%   \sem{e}               [[e]]               (denotation)
%   \assign, \seqc        :=  and the sequencing semicolon
%   \impl                 ==> as a math arrow (LaTeX's own \iff gives <==>)
% ---------------------------------------------------------------------------
\newcommand{\hoare}[3]{\{\,#1\,\}\;#2\;\{\,#3\,\}}
\newcommand{\thoare}[3]{[\,#1\,]\;#2\;[\,#3\,]}
\newcommand{\provable}[3]{\vdash \hoare{#1}{#2}{#3}}
\newcommand{\WP}[2]{\mathrm{wp}(#1,\,#2)}
\newcommand{\WLP}[2]{\mathrm{wlp}(#1,\,#2)}
\newcommand{\VC}[2]{\mathrm{VC}(#1,\,#2)}
\newcommand{\subst}[3]{#1[#3/#2]}
\newcommand{\sem}[1]{\llbracket #1 \rrbracket}
\newcommand{\assign}{\mathrel{:=}}
\newcommand{\seqc}{\mathbin{;}\,}
\newcommand{\impl}{\Longrightarrow}

% ---------------------------------------------------------------------------
% Judgments.  \judgment{\Gamma}{e}{\tau}  ==>  Gamma |- e : tau
% ---------------------------------------------------------------------------
\newcommand{\judgment}[3]{#1 \vdash #2 : #3}
\newcommand{\judg}[3]{\judgment{#1}{#2}{#3}}          % short alias
\newcommand{\emptyjudgment}[2]{\cdot \vdash #1 : #2}  % empty context
\newcommand{\ctx}{\Gamma}
\newcommand{\turnstile}{\vdash}

% ---------------------------------------------------------------------------
% Theorem environments.
% ---------------------------------------------------------------------------
\theoremstyle{plain}
\newtheorem{theorem}{Theorem}[section]
\newtheorem{lemma}[theorem]{Lemma}
\newtheorem{corollary}[theorem]{Corollary}
\newtheorem{proposition}[theorem]{Proposition}
\theoremstyle{definition}
\newtheorem{definition}[theorem]{Definition}
% Written-track problems are numbered B1, B2, ... on both tracks (P<n> is a
% predict item). Label each one \label{prob:B<n>} and its answer \label{ans:B<n>}.
\newtheorem{problem}{Problem}
\renewcommand{\theproblem}{B\arabic{problem}}
\theoremstyle{remark}
\newtheorem*{remark}{Remark}
\newenvironment{answer}
  {\par\medskip\noindent\textbf{Answer.}\begin{quotation}\noindent}
  {\end{quotation}\medskip}

% ---------------------------------------------------------------------------
% Grader margin notes. Rendered form of the `% GRADER[...][...][...]:` comment
% convention (see template_answers.tex). Use inline where a visible note helps:
%   \gradernote{B2/induction-soundness/2: the invariant is not preserved by the last assignment}
% ---------------------------------------------------------------------------
\definecolor{gradercolor}{RGB}{160, 30, 30}
\newcommand{\gradernote}[1]{\textcolor{gradercolor}{\small\sffamily [GRADER: #1]}}

% ---------------------------------------------------------------------------
% Code listings: Dafny, Verus (Rust inside verus!{...}), Python.
% ---------------------------------------------------------------------------
\definecolor{codekeyword}{RGB}{0, 60, 130}
\definecolor{codecomment}{RGB}{90, 110, 90}
\definecolor{codestring}{RGB}{130, 60, 10}
\definecolor{codehole}{RGB}{180, 0, 0}

\lstdefinelanguage{dafny}{
  sensitive=true,
  morekeywords={
    abstract, allocated, as, assert, assume, bool, break, by, calc, case, char,
    class, codatatype, const, constructor, continue, datatype, decreases, else,
    ensures, exists, expect, export, extends, false, forall, fresh, function,
    ghost, if, import, in, include, int, invariant, is, iterator, label, lemma,
    map, match, method, modifies, modify, module, multiset, nat, new, newtype,
    null, old, opaque, opened, predicate, print, provides, reads, real, refines,
    requires, return, returns, reveal, reveals, seq, set, static, string, then,
    this, trait, true, twostate, type, unchanged, var, while, witness},
  emph={assume},
  emphstyle=\color{codehole}\bfseries,
  comment=[l]{//},
  morecomment=[s]{/*}{*/},
  morestring=[b]",
}

\lstdefinelanguage{verus}{
  sensitive=true,
  morekeywords={
    as, break, const, continue, crate, dyn, else, enum, false, fn, for, if,
    impl, in, let, loop, match, mod, move, mut, pub, ref, return, self, Self,
    static, struct, super, trait, true, type, unsafe, use, where, while,
    spec, proof, exec, open, closed, tracked, ghost, requires, ensures,
    recommends, invariant, decreases, assert, forall, exists, choose, by,
    broadcast, uninterp, axiom, returns, int, nat, Seq, Set, Map, Ghost, Tracked},
  emph={assume, admit},
  emphstyle=\color{codehole}\bfseries,
  comment=[l]{//},
  morecomment=[s]{/*}{*/},
  morestring=[b]",
}

\lstset{
  basicstyle=\ttfamily\small,
  keywordstyle=\color{codekeyword}\bfseries,
  commentstyle=\color{codecomment}\itshape,
  stringstyle=\color{codestring},
  showstringspaces=false,
  columns=fullflexible,
  keepspaces=true,
  breaklines=true,
  frame=leftline,
  framerule=0.6pt,
  xleftmargin=1em,
  aboveskip=0.8em,
  belowskip=0.8em,
  literate=
    {→}{{$\to$}}1        {←}{{$\leftarrow$}}1  {⇒}{{$\Rightarrow$}}1
    {⇔}{{$\Leftrightarrow$}}1                  {↦}{{$\mapsto$}}1
    {∀}{{$\forall$}}1    {∃}{{$\exists$}}1     {λ}{{$\lambda$}}1
    {∧}{{$\wedge$}}1     {∨}{{$\vee$}}1        {¬}{{$\neg$}}1
    {≠}{{$\neq$}}1       {≤}{{$\leq$}}1        {≥}{{$\geq$}}1
    {∈}{{$\in$}}1        {∉}{{$\notin$}}1      {⊆}{{$\subseteq$}}1
    {×}{{$\times$}}1     {·}{{$\cdot$}}1       {∅}{{$\varnothing$}}1
    {ℕ}{{$\mathbb{N}$}}1 {ℤ}{{$\mathbb{Z}$}}1
    {⊢}{{$\vdash$}}1     {⟦}{{$\llbracket$}}1  {⟧}{{$\rrbracket$}}1,
}

% Code environments and inline code.
\lstnewenvironment{dafnycode}{\lstset{language=dafny}}{}
\lstnewenvironment{veruscode}{\lstset{language=verus}}{}
\lstnewenvironment{pycode}{\lstset{language=Python}}{}
\newcommand{\dfy}[1]{\lstinline[language=dafny]{#1}}
\newcommand{\vrs}[1]{\lstinline[language=verus]{#1}}
\newcommand{\py}[1]{\lstinline[language=Python]{#1}}
          % the Makefile always runs latexmk from paper/
%   \begin{document}
%   ...
%   \end{document}
%
% Engine: pdflatex. Dafny and Verus sources are mostly ASCII; the `literate`
% tables below map the common Unicode symbols to math glyphs inside listings.
% For other Unicode, extend a table, or build that file with lualatex
% (`latexmk -lualatex <file>`) after replacing the inputenc/fontenc lines with
% \usepackage{fontspec} and a monospace font with good coverage.
%
% Contents: program-logic macros (Hoare triples, wp, VCs, substitution,
% semantics brackets, typing judgments), inference rules (mathpartir),
% theorem and problem environments, the answer environment, the grader
% margin-note macro, and listings languages `dafny` and `verus` (Python uses
% listings' built-in `Python`), with inline forms \dfy{...}, \vrs{...}, \py{...}.
\usepackage[utf8]{inputenc}
\usepackage[T1]{fontenc}

\usepackage{amsmath}
\usepackage{amssymb}
\usepackage{amsthm}
\usepackage{stmaryrd}         % \llbracket \rrbracket for semantics brackets
\usepackage{mathpartir}       % inference rules: \inferrule, the mathpar environment
\usepackage{listings}
\usepackage{xcolor}
\usepackage[margin=1.1in]{geometry}

% ---------------------------------------------------------------------------
% Program logic.
%   \hoare{P}{c}{Q}       {P} c {Q}           (partial correctness)
%   \thoare{P}{c}{Q}      [P] c [Q]           (total correctness)
%   \provable{P}{c}{Q}    |- {P} c {Q}        (derivable in the proof system)
%   \WP{c}{Q}             wp(c, Q)            (\wp is LaTeX's Weierstrass p, so
%   \WLP{c}{Q}            wlp(c, Q)            the macro names are upper case)
%   \VC{c}{Q}             VC(c, Q)
%   \subst{Q}{x}{e}       Q[e/x]              (e replaces x in Q)
%   \sem{e}               [[e]]               (denotation)
%   \assign, \seqc        :=  and the sequencing semicolon
%   \impl                 ==> as a math arrow (LaTeX's own \iff gives <==>)
% ---------------------------------------------------------------------------
\newcommand{\hoare}[3]{\{\,#1\,\}\;#2\;\{\,#3\,\}}
\newcommand{\thoare}[3]{[\,#1\,]\;#2\;[\,#3\,]}
\newcommand{\provable}[3]{\vdash \hoare{#1}{#2}{#3}}
\newcommand{\WP}[2]{\mathrm{wp}(#1,\,#2)}
\newcommand{\WLP}[2]{\mathrm{wlp}(#1,\,#2)}
\newcommand{\VC}[2]{\mathrm{VC}(#1,\,#2)}
\newcommand{\subst}[3]{#1[#3/#2]}
\newcommand{\sem}[1]{\llbracket #1 \rrbracket}
\newcommand{\assign}{\mathrel{:=}}
\newcommand{\seqc}{\mathbin{;}\,}
\newcommand{\impl}{\Longrightarrow}

% ---------------------------------------------------------------------------
% Judgments.  \judgment{\Gamma}{e}{\tau}  ==>  Gamma |- e : tau
% ---------------------------------------------------------------------------
\newcommand{\judgment}[3]{#1 \vdash #2 : #3}
\newcommand{\judg}[3]{\judgment{#1}{#2}{#3}}          % short alias
\newcommand{\emptyjudgment}[2]{\cdot \vdash #1 : #2}  % empty context
\newcommand{\ctx}{\Gamma}
\newcommand{\turnstile}{\vdash}

% ---------------------------------------------------------------------------
% Theorem environments.
% ---------------------------------------------------------------------------
\theoremstyle{plain}
\newtheorem{theorem}{Theorem}[section]
\newtheorem{lemma}[theorem]{Lemma}
\newtheorem{corollary}[theorem]{Corollary}
\newtheorem{proposition}[theorem]{Proposition}
\theoremstyle{definition}
\newtheorem{definition}[theorem]{Definition}
% Written-track problems are numbered B1, B2, ... on both tracks (P<n> is a
% predict item). Label each one \label{prob:B<n>} and its answer \label{ans:B<n>}.
\newtheorem{problem}{Problem}
\renewcommand{\theproblem}{B\arabic{problem}}
\theoremstyle{remark}
\newtheorem*{remark}{Remark}
\newenvironment{answer}
  {\par\medskip\noindent\textbf{Answer.}\begin{quotation}\noindent}
  {\end{quotation}\medskip}

% ---------------------------------------------------------------------------
% Grader margin notes. Rendered form of the `% GRADER[...][...][...]:` comment
% convention (see template_answers.tex). Use inline where a visible note helps:
%   \gradernote{B2/induction-soundness/2: the invariant is not preserved by the last assignment}
% ---------------------------------------------------------------------------
\definecolor{gradercolor}{RGB}{160, 30, 30}
\newcommand{\gradernote}[1]{\textcolor{gradercolor}{\small\sffamily [GRADER: #1]}}

% ---------------------------------------------------------------------------
% Code listings: Dafny, Verus (Rust inside verus!{...}), Python.
% ---------------------------------------------------------------------------
\definecolor{codekeyword}{RGB}{0, 60, 130}
\definecolor{codecomment}{RGB}{90, 110, 90}
\definecolor{codestring}{RGB}{130, 60, 10}
\definecolor{codehole}{RGB}{180, 0, 0}

\lstdefinelanguage{dafny}{
  sensitive=true,
  morekeywords={
    abstract, allocated, as, assert, assume, bool, break, by, calc, case, char,
    class, codatatype, const, constructor, continue, datatype, decreases, else,
    ensures, exists, expect, export, extends, false, forall, fresh, function,
    ghost, if, import, in, include, int, invariant, is, iterator, label, lemma,
    map, match, method, modifies, modify, module, multiset, nat, new, newtype,
    null, old, opaque, opened, predicate, print, provides, reads, real, refines,
    requires, return, returns, reveal, reveals, seq, set, static, string, then,
    this, trait, true, twostate, type, unchanged, var, while, witness},
  emph={assume},
  emphstyle=\color{codehole}\bfseries,
  comment=[l]{//},
  morecomment=[s]{/*}{*/},
  morestring=[b]",
}

\lstdefinelanguage{verus}{
  sensitive=true,
  morekeywords={
    as, break, const, continue, crate, dyn, else, enum, false, fn, for, if,
    impl, in, let, loop, match, mod, move, mut, pub, ref, return, self, Self,
    static, struct, super, trait, true, type, unsafe, use, where, while,
    spec, proof, exec, open, closed, tracked, ghost, requires, ensures,
    recommends, invariant, decreases, assert, forall, exists, choose, by,
    broadcast, uninterp, axiom, returns, int, nat, Seq, Set, Map, Ghost, Tracked},
  emph={assume, admit},
  emphstyle=\color{codehole}\bfseries,
  comment=[l]{//},
  morecomment=[s]{/*}{*/},
  morestring=[b]",
}

\lstset{
  basicstyle=\ttfamily\small,
  keywordstyle=\color{codekeyword}\bfseries,
  commentstyle=\color{codecomment}\itshape,
  stringstyle=\color{codestring},
  showstringspaces=false,
  columns=fullflexible,
  keepspaces=true,
  breaklines=true,
  frame=leftline,
  framerule=0.6pt,
  xleftmargin=1em,
  aboveskip=0.8em,
  belowskip=0.8em,
  literate=
    {→}{{$\to$}}1        {←}{{$\leftarrow$}}1  {⇒}{{$\Rightarrow$}}1
    {⇔}{{$\Leftrightarrow$}}1                  {↦}{{$\mapsto$}}1
    {∀}{{$\forall$}}1    {∃}{{$\exists$}}1     {λ}{{$\lambda$}}1
    {∧}{{$\wedge$}}1     {∨}{{$\vee$}}1        {¬}{{$\neg$}}1
    {≠}{{$\neq$}}1       {≤}{{$\leq$}}1        {≥}{{$\geq$}}1
    {∈}{{$\in$}}1        {∉}{{$\notin$}}1      {⊆}{{$\subseteq$}}1
    {×}{{$\times$}}1     {·}{{$\cdot$}}1       {∅}{{$\varnothing$}}1
    {ℕ}{{$\mathbb{N}$}}1 {ℤ}{{$\mathbb{Z}$}}1
    {⊢}{{$\vdash$}}1     {⟦}{{$\llbracket$}}1  {⟧}{{$\rrbracket$}}1,
}

% Code environments and inline code.
\lstnewenvironment{dafnycode}{\lstset{language=dafny}}{}
\lstnewenvironment{veruscode}{\lstset{language=verus}}{}
\lstnewenvironment{pycode}{\lstset{language=Python}}{}
\newcommand{\dfy}[1]{\lstinline[language=dafny]{#1}}
\newcommand{\vrs}[1]{\lstinline[language=verus]{#1}}
\newcommand{\py}[1]{\lstinline[language=Python]{#1}}
          % the Makefile always runs latexmk from paper/
%   \begin{document}
%   ...
%   \end{document}
%
% Engine: pdflatex. Dafny and Verus sources are mostly ASCII; the `literate`
% tables below map the common Unicode symbols to math glyphs inside listings.
% For other Unicode, extend a table, or build that file with lualatex
% (`latexmk -lualatex <file>`) after replacing the inputenc/fontenc lines with
% \usepackage{fontspec} and a monospace font with good coverage.
%
% Contents: program-logic macros (Hoare triples, wp, VCs, substitution,
% semantics brackets, typing judgments), inference rules (mathpartir),
% theorem and problem environments, the answer environment, the grader
% margin-note macro, and listings languages `dafny` and `verus` (Python uses
% listings' built-in `Python`), with inline forms \dfy{...}, \vrs{...}, \py{...}.
\usepackage[utf8]{inputenc}
\usepackage[T1]{fontenc}

\usepackage{amsmath}
\usepackage{amssymb}
\usepackage{amsthm}
\usepackage{stmaryrd}         % \llbracket \rrbracket for semantics brackets
\usepackage{mathpartir}       % inference rules: \inferrule, the mathpar environment
\usepackage{listings}
\usepackage{xcolor}
\usepackage[margin=1.1in]{geometry}

% ---------------------------------------------------------------------------
% Program logic.
%   \hoare{P}{c}{Q}       {P} c {Q}           (partial correctness)
%   \thoare{P}{c}{Q}      [P] c [Q]           (total correctness)
%   \provable{P}{c}{Q}    |- {P} c {Q}        (derivable in the proof system)
%   \WP{c}{Q}             wp(c, Q)            (\wp is LaTeX's Weierstrass p, so
%   \WLP{c}{Q}            wlp(c, Q)            the macro names are upper case)
%   \VC{c}{Q}             VC(c, Q)
%   \subst{Q}{x}{e}       Q[e/x]              (e replaces x in Q)
%   \sem{e}               [[e]]               (denotation)
%   \assign, \seqc        :=  and the sequencing semicolon
%   \impl                 ==> as a math arrow (LaTeX's own \iff gives <==>)
% ---------------------------------------------------------------------------
\newcommand{\hoare}[3]{\{\,#1\,\}\;#2\;\{\,#3\,\}}
\newcommand{\thoare}[3]{[\,#1\,]\;#2\;[\,#3\,]}
\newcommand{\provable}[3]{\vdash \hoare{#1}{#2}{#3}}
\newcommand{\WP}[2]{\mathrm{wp}(#1,\,#2)}
\newcommand{\WLP}[2]{\mathrm{wlp}(#1,\,#2)}
\newcommand{\VC}[2]{\mathrm{VC}(#1,\,#2)}
\newcommand{\subst}[3]{#1[#3/#2]}
\newcommand{\sem}[1]{\llbracket #1 \rrbracket}
\newcommand{\assign}{\mathrel{:=}}
\newcommand{\seqc}{\mathbin{;}\,}
\newcommand{\impl}{\Longrightarrow}

% ---------------------------------------------------------------------------
% Judgments.  \judgment{\Gamma}{e}{\tau}  ==>  Gamma |- e : tau
% ---------------------------------------------------------------------------
\newcommand{\judgment}[3]{#1 \vdash #2 : #3}
\newcommand{\judg}[3]{\judgment{#1}{#2}{#3}}          % short alias
\newcommand{\emptyjudgment}[2]{\cdot \vdash #1 : #2}  % empty context
\newcommand{\ctx}{\Gamma}
\newcommand{\turnstile}{\vdash}

% ---------------------------------------------------------------------------
% Theorem environments.
% ---------------------------------------------------------------------------
\theoremstyle{plain}
\newtheorem{theorem}{Theorem}[section]
\newtheorem{lemma}[theorem]{Lemma}
\newtheorem{corollary}[theorem]{Corollary}
\newtheorem{proposition}[theorem]{Proposition}
\theoremstyle{definition}
\newtheorem{definition}[theorem]{Definition}
% Written-track problems are numbered B1, B2, ... on both tracks (P<n> is a
% predict item). Label each one \label{prob:B<n>} and its answer \label{ans:B<n>}.
\newtheorem{problem}{Problem}
\renewcommand{\theproblem}{B\arabic{problem}}
\theoremstyle{remark}
\newtheorem*{remark}{Remark}
\newenvironment{answer}
  {\par\medskip\noindent\textbf{Answer.}\begin{quotation}\noindent}
  {\end{quotation}\medskip}

% ---------------------------------------------------------------------------
% Grader margin notes. Rendered form of the `% GRADER[...][...][...]:` comment
% convention (see template_answers.tex). Use inline where a visible note helps:
%   \gradernote{B2/induction-soundness/2: the invariant is not preserved by the last assignment}
% ---------------------------------------------------------------------------
\definecolor{gradercolor}{RGB}{160, 30, 30}
\newcommand{\gradernote}[1]{\textcolor{gradercolor}{\small\sffamily [GRADER: #1]}}

% ---------------------------------------------------------------------------
% Code listings: Dafny, Verus (Rust inside verus!{...}), Python.
% ---------------------------------------------------------------------------
\definecolor{codekeyword}{RGB}{0, 60, 130}
\definecolor{codecomment}{RGB}{90, 110, 90}
\definecolor{codestring}{RGB}{130, 60, 10}
\definecolor{codehole}{RGB}{180, 0, 0}

\lstdefinelanguage{dafny}{
  sensitive=true,
  morekeywords={
    abstract, allocated, as, assert, assume, bool, break, by, calc, case, char,
    class, codatatype, const, constructor, continue, datatype, decreases, else,
    ensures, exists, expect, export, extends, false, forall, fresh, function,
    ghost, if, import, in, include, int, invariant, is, iterator, label, lemma,
    map, match, method, modifies, modify, module, multiset, nat, new, newtype,
    null, old, opaque, opened, predicate, print, provides, reads, real, refines,
    requires, return, returns, reveal, reveals, seq, set, static, string, then,
    this, trait, true, twostate, type, unchanged, var, while, witness},
  emph={assume},
  emphstyle=\color{codehole}\bfseries,
  comment=[l]{//},
  morecomment=[s]{/*}{*/},
  morestring=[b]",
}

\lstdefinelanguage{verus}{
  sensitive=true,
  morekeywords={
    as, break, const, continue, crate, dyn, else, enum, false, fn, for, if,
    impl, in, let, loop, match, mod, move, mut, pub, ref, return, self, Self,
    static, struct, super, trait, true, type, unsafe, use, where, while,
    spec, proof, exec, open, closed, tracked, ghost, requires, ensures,
    recommends, invariant, decreases, assert, forall, exists, choose, by,
    broadcast, uninterp, axiom, returns, int, nat, Seq, Set, Map, Ghost, Tracked},
  emph={assume, admit},
  emphstyle=\color{codehole}\bfseries,
  comment=[l]{//},
  morecomment=[s]{/*}{*/},
  morestring=[b]",
}

\lstset{
  basicstyle=\ttfamily\small,
  keywordstyle=\color{codekeyword}\bfseries,
  commentstyle=\color{codecomment}\itshape,
  stringstyle=\color{codestring},
  showstringspaces=false,
  columns=fullflexible,
  keepspaces=true,
  breaklines=true,
  frame=leftline,
  framerule=0.6pt,
  xleftmargin=1em,
  aboveskip=0.8em,
  belowskip=0.8em,
  literate=
    {→}{{$\to$}}1        {←}{{$\leftarrow$}}1  {⇒}{{$\Rightarrow$}}1
    {⇔}{{$\Leftrightarrow$}}1                  {↦}{{$\mapsto$}}1
    {∀}{{$\forall$}}1    {∃}{{$\exists$}}1     {λ}{{$\lambda$}}1
    {∧}{{$\wedge$}}1     {∨}{{$\vee$}}1        {¬}{{$\neg$}}1
    {≠}{{$\neq$}}1       {≤}{{$\leq$}}1        {≥}{{$\geq$}}1
    {∈}{{$\in$}}1        {∉}{{$\notin$}}1      {⊆}{{$\subseteq$}}1
    {×}{{$\times$}}1     {·}{{$\cdot$}}1       {∅}{{$\varnothing$}}1
    {ℕ}{{$\mathbb{N}$}}1 {ℤ}{{$\mathbb{Z}$}}1
    {⊢}{{$\vdash$}}1     {⟦}{{$\llbracket$}}1  {⟧}{{$\rrbracket$}}1,
}

% Code environments and inline code.
\lstnewenvironment{dafnycode}{\lstset{language=dafny}}{}
\lstnewenvironment{veruscode}{\lstset{language=verus}}{}
\lstnewenvironment{pycode}{\lstset{language=Python}}{}
\newcommand{\dfy}[1]{\lstinline[language=dafny]{#1}}
\newcommand{\vrs}[1]{\lstinline[language=verus]{#1}}
\newcommand{\py}[1]{\lstinline[language=Python]{#1}}


\title{Q15 --- Practice II-C + Lab 2 \\ \large Written problems}
\author{dafny-verus-learn-0}
\date{}

\begin{document}
\maketitle
\sloppy

\section*{Instructions}
Answer in \texttt{Q15\_answers.tex}, one section per problem; keep its
\verb+\label{ans:B<n>}+ anchors. \texttt{make -C paper paper} must stay
green: a non-compiling answer set is the written track's red build.

Order: the three cards (B3--B5) first, 15 minutes each, each reread in your own
\texttt{cheatsheets.md}; then the six prose specifications and your oracles;
then B1 (40 minutes), before any Lab 2 run. Run \texttt{/grade-paper Q15} in a
fresh session on B1 and the cards: the runs unlock only on its unlocked
gradebook row. B2 comes after the runs and is graded only by
\texttt{/grade-paper Q15 --report}. Compare with \texttt{paper/solutions/}
only as the module's last step.

Lab 2 asks three things about specifications written from prose. Do your
oracles (\texttt{harness/q/fvh\_q/oracles.py}) separate faithful, weak,
vacuous and over-strong specifications, judged by their agreement with the
hidden verdicts? Does reference-free oracle feedback make an agent's
specifications more faithful? And how often does an agent that implements and
proves against its own specification fail when the reference re-judges it
(vericoding)? The Qwen arms are S0 (a specification from the prose, four
independent samples, the first that verifies with its witness kept), S1 (the
same four calls, with reference-free feedback after each sample) and V (the
agent writes the specification, then implements and proves against it, with
verifier feedback, eight calls); a Claude arm runs S1 on a seeded,
language-stratified subset of 20 tasks. S0 and S1 are matched at four calls
per task; V is a separate measurement. The tasks: a local stratum
(\texttt{bench/local/q/lab2/STRATUM.toml}, ten Dafny and ten Verus prose
tasks) and two public strata of five tasks each (QPLAN §25.6). The hidden
reference never enters a loop: \texttt{implication} and \texttt{classify}
run only after a batch, in the grader container.

% ---------------------------------------------------------------------------
\begin{problem}[pre-registration]\label{prob:B1}
Pre-register Lab 2 in the PREREG section of your answers file, on every field
of the template there, with your own numbers: the claim (oracle separation;
the feedback effect S1 $-$ S0 at matched calls, one- or two-sided; the
vericoding re-judged failure rate, Dafny against Verus, reported on its own
and never compared with S0 or S1); the feedback signals S1 receives, each
named and each shown to need no hidden input; the metric and estimator (the
oracle verdict against the hidden reference, agreement with the hidden
verdicts as a confusion matrix with Cohen's $\kappa$, the
verified-but-unfaithful rate with Wilson intervals, and the paired estimator
of S1 $-$ S0); the arms and policies with their prompt hashes and the Claude
subset's seed and stratification; the tasks and strata; the budget
arithmetic per arm, with the sentinel calls and the throughput assumption and
its source, against the registry's \texttt{budget}; the generation settings;
the cost per arm with its truncation rate, and the fallback if the per-arm
totals exceed the ceiling; the MDE, computed with
\texttt{fvh.stats.mde\_two\_proportions} at your $n$, baseline, pairing and
$\alpha$; the seeds and repeats; the oracles' version (the sha256 of
\texttt{oracles.py} at the run); the cheat surface and gate; the exclusions;
the stopping rule; your 80\% prediction intervals; the contamination
controls, with the card contamination audit of the local tasks
(authorship after the models' release, a novelty review, perturbation
probes); the batch plan path and its final \texttt{plan-sha256:} binding;
and the disclosure. The runs unlock only when
\texttt{/grade-paper Q15} scores every method dimension at least 3, and the
harness refuses a run whose pre-registration hash is not on that unlocked
gradebook row.
\end{problem}

% ---------------------------------------------------------------------------
\begin{problem}[lab-report]\label{prob:B2}
After the runs, report Lab 2 against its pre-registration: every table
regenerated by your \emph{Reproduce} command from the logs in
\texttt{research/runs/Q15/}, the confusion matrix and $\kappa$ among them;
every deviation; the gate's verdict on every claimed success, and every
claimed faithful verdict reproducible in the grader container from a clean
checkout; S1 $-$ S0 with its interval; the vericoding re-judged failure
rate per language, with its Wilson interval, on its own; the reward audit
(the planted specifications that would score well on S1's feedback and still
be unfaithful, and the oracle that catches each) and the bench audit (your
oracles on the course's own \texttt{bench/local/q} cards, each finding
named); your prediction intervals scored; contamination; and the disclosure
block. It is graded with \texttt{/grade-paper Q15 --report}, on method and
reproducibility, never on the numbers.
\end{problem}

% ---------------------------------------------------------------------------
\begin{problem}[paper-card]\label{prob:B3}
\textbf{Paper card 04.6, Re:Form: Reducing Human Annotations in Scalable Formal Software Verification with RL in LLMs} (15 minutes). Fill the card on its eight
fields: task shape, oracle, what is trusted, cheat surface, metric, budget,
variance, artifact. Each field is graded against the pinned range of your own
\texttt{cheatsheets.md}: lines 1618--1649, the entry from its heading to its
\emph{Why it matters} paragraph, in the copy the command below names. The pin
is the sha256 of exactly those lines, computed with the first line below; it
must print the second.
{\small
\begin{verbatim}
sed -n '1618,1649p' ~/code/research/agentic-formal-veri/cheatsheets.md | sha256sum
0652aa5774df67acf6924a7912ace1d9bc9f1910dbc9aa28ddd444cf561329b3  -
\end{verbatim}
}
Write each field in your own words, a sentence or two. Where the range says
nothing about a field, write that, and what you would need to know. A hash
that no longer matches means the card file changed: the grader refuses the
card until the pin is renewed. Then one more line: which of Lab 2's oracles
(\texttt{implication}, \texttt{kill\_rate}, \texttt{vacuity}) is the paper's
quality signal closest to, and could S1 receive it, that is, does it need the
reference?
\end{problem}

% ---------------------------------------------------------------------------
\begin{problem}[paper-card]\label{prob:B4}
\textbf{Paper card 04.7, SpecRL: Reinforcement Learning with Test-Based Completeness Rewards for Formal Specification Synthesis} (15 minutes). Fill the card on its eight
fields: task shape, oracle, what is trusted, cheat surface, metric, budget,
variance, artifact. Each field is graded against the pinned range of your own
\texttt{cheatsheets.md}: lines 1653--1677, the entry from its heading to its
\emph{Why it matters} paragraph, in the copy the command below names. The pin
is the sha256 of exactly those lines, computed with the first line below; it
must print the second.
{\small
\begin{verbatim}
sed -n '1653,1677p' ~/code/research/agentic-formal-veri/cheatsheets.md | sha256sum
2ee2bc96c9712432918e5b28957c63b7df3e7a72811311e9e82390e6c13c9202  -
\end{verbatim}
}
Write each field in your own words, a sentence or two. Where the range says
nothing about a field, write that, and what you would need to know. A hash
that no longer matches means the card file changed: the grader refuses the
card until the pin is renewed. Then one more line: which of Lab 2's oracles
(\texttt{implication}, \texttt{kill\_rate}, \texttt{vacuity}) is the paper's
quality signal closest to, and could S1 receive it, that is, does it need the
reference?
\end{problem}

% ---------------------------------------------------------------------------
\begin{problem}[paper-card]\label{prob:B5}
\textbf{Paper card 03.8, Can Large Language Models Transform Natural Language Intent into Formal Method Postconditions?} (15 minutes). Fill the card on its eight
fields: task shape, oracle, what is trusted, cheat surface, metric, budget,
variance, artifact. Each field is graded against the pinned range of your own
\texttt{cheatsheets.md}: lines 1257--1282, the entry from its heading to its
\emph{Why it matters} paragraph, in the copy the command below names. The pin
is the sha256 of exactly those lines, computed with the first line below; it
must print the second.
{\small
\begin{verbatim}
sed -n '1257,1282p' ~/code/research/agentic-formal-veri/cheatsheets.md | sha256sum
837a45587f1cad82ced01ce9f6e2b15ac6a5c3a9ad914e0fff23e11c9812dde4  -
\end{verbatim}
}
Write each field in your own words, a sentence or two. Where the range says
nothing about a field, write that, and what you would need to know. A hash
that no longer matches means the card file changed: the grader refuses the
card until the pin is renewed. Then one more line: which of Lab 2's oracles
(\texttt{implication}, \texttt{kill\_rate}, \texttt{vacuity}) is the paper's
quality signal closest to, and could S1 receive it, that is, does it need the
reference?
\end{problem}

\end{document}
